\input{../tex-inputs/latex-leseansicht-vorspann}

\section[Arthur und Olga Schnitzler an Gustav Schwarzkopf, 29. 12. 1904]{L04384 Arthur und Olga Schnitzler an Gustav Schwarzkopf, 29. 12. 1904}
\nopagebreak\mylabel{L04384v}
\rehead{ }\normalsize\beginnumbering\briefempfaengerindex{Schwarzkopf, Gustav@\textsc{Schwarzkopf, Gustav}!zzzSchnitzler, Olga@\emph{von Olga Schnitzler}!1904-12-291@{29. 12. 1904}|(be}\briefempfaengerindex{Schwarzkopf, Gustav@\textsc{Schwarzkopf, Gustav}!zzzSchnitzler, Arthur@\emph{von Arthur Schnitzler}!1904-12-291@{29. 12. 1904}|(be}
\toendnotes[C]{\smallbreak\pagebreak[2]}
\correspDesc{Versand  durch Arthur Schnitzler, Olga Schnitzler am 29. 12. 1904 in St. Gilgen
\newline{}Übermittlung  am 29. 12. 1904 in St. Gilgen
\newline{}Zustellung  am 30. 12. 1904 in Wien
\newline{}Erhalt  durch Gustav Schwarzkopf am 30. 12. 1904 in Wien}\toendnotes[C]{\smallbreak}
\Standort{DLA, A:Schnitzler, HS.1985.1.1897.}
\physDesc{Bildpostkarte, 176 Zeichen
\newline{}Handschrift Arthur Schnitzler: Bleistift, lateinische Kurrent
\newline{}Handschrift Olga Schnitzler: Bleistift, lateinische Kurrent
\newline{}Versand: 1) Stempel: »\nobreak{}\oindex{St. Gilgen@\textbf{St. Gilgen}|pwk}St. Gilgen, 29 {[}XII{]} 04\nobreak{}«.   2) Stempel: »\nobreak{}\oindex{I., Innere Stadt@\textbf{I., Innere Stadt}|pwk}{[}Wien{]}{ }\textcolor{gray}{1}/1
                                       1, \textcolor{gray}{3}0. 12. \textcolor{gray}{04}, 9–\textcolor{gray}{10}V., \textcolor{gray}{Best}ellt\nobreak{}«. }\toendnotes[C]{\smallbreak}\pstart{}{\pb}Herrn Gustav Schwarzkopf\pend{}\pstart{}Wien \strikeout{IX} I.\oindex{I., Innere Stadt@\textbf{I., Innere Stadt}|pw}\pend{}\pstart{}Tiefer Graben 23\oindex{Wien@\textbf{Wien}!I., Innere Stadt@\textbf{I., Innere Stadt}!Tiefer Graben 23@\textbf{Tiefer Graben 23}|pw}.\pend{}{\bigskip}
\pstart
           \centering{}{\pb}\textcolor{gray}{\textbf{St. Gilgen}}\oindex{St. Gilgen@\textbf{St. Gilgen}|pw}\pend
           
\pstart
           \centering{}\textcolor{gray}{\textbf{Salzburg}}\oindex{Salzburg [Land]@\textbf{Salzburg [Land]}|pw}\pend
           \vspace{1em}
\pstart
           \raggedleft{}{\pb}Lueg\oindex{Lueg@\textbf{Lueg}|pw}{ }\label{T_L04371-1v}\edtext{29. 12. 904}{\lemma{\textnormal{\emph{29. 12. 904}}}\Cendnote{\textnormal{am linken oberen Bildrand, verkehrt zum Motiv}}}\label{T_L04371-1}\pend
           \vspace{0.5em}
\pstart
           Herzliche Grüße und auf baldiges  Wiedersehen!\pend
           
\pstart
           Ihr{\\[\baselineskip]}\spacefill\mbox{Arthur S.}\pend
           \leftskip=0em{}\selectlanguage{ngerman}\vspace{1em}
\pstart
           \noindent{}{[}hs. O. Schnitzler:{]} Gut Neujahr u. viele Grüsse! auch an Herrn D\textsuperscript{r}{ }M.\pwindex{Schwarzkopf, Max 12.\,6.\,1857 Wien – 14.\,4.\,1928 ebd.@\textsc{Schwarzkopf, Max} (12.\,6.\,1857 Wien – 14.\,4.\,1928 ebd.)|pw}\pend
           \pstart \spacefill\mbox{O. S.}\pend{}\selectlanguage{ngerman}\endnumbering\briefempfaengerindex{Schwarzkopf, Gustav@\textsc{Schwarzkopf, Gustav}!zzzSchnitzler, Olga@\emph{von Olga Schnitzler}!1904-12-291@{29. 12. 1904}|)be}\briefempfaengerindex{Schwarzkopf, Gustav@\textsc{Schwarzkopf, Gustav}!zzzSchnitzler, Arthur@\emph{von Arthur Schnitzler}!1904-12-291@{29. 12. 1904}|)be}\mylabel{L04384h}
\begin{anhang}
\end{anhang}\newcommand{\dateiname}{L04384}\newcommand{\titel}{Arthur und Olga Schnitzler an Gustav Schwarzkopf, 29. 12. 1904}\newcommand{\editorInnen}{Herausgegeben von Jahnke, SelmaMüller, Martin Anton}\input{../tex-inputs/latex-leseansicht-abspann}
